\PassOptionsToPackage{hidelinks}{hyperref}
\documentclass[
  aps,
  prd,
  reprint,
  superscriptaddress,
  nofootinbib,
  floatfix,
  longbibliography
]{revtex4-2}

\usepackage[T1]{fontenc}
\usepackage[utf8]{inputenc}
\usepackage{amsmath,amssymb,amsthm,mathtools,bm}
\usepackage{booktabs,array,tabularx}
\usepackage{graphicx}
\usepackage{microtype}
\IfFileExists{orcidlink.sty}{\usepackage{orcidlink}}{\newcommand{\orcidlink}[1]{}}
\usepackage{hyperref}

\hypersetup{
  pdftitle={Dark-sector residuals in compact S3 FLRW cosmology},
  pdfauthor={Andrei Preece and Boris Batenin}
}

\emergencystretch=3em
\allowdisplaybreaks

\theoremstyle{plain}
\newtheorem{theorem}{Theorem}[section]
\newtheorem{lemma}[theorem]{Lemma}
\newtheorem{proposition}[theorem]{Proposition}
\newtheorem{corollary}[theorem]{Corollary}
\theoremstyle{definition}
\newtheorem{definition}[theorem]{Definition}
\theoremstyle{remark}
\newtheorem{remark}[theorem]{Remark}
\newtheorem{nonclaim}[theorem]{Nonclaim}

\newcommand{\dd}{\mathrm d}
\newcommand{\RR}{\mathbb R}
\newcommand{\M}{\mathcal M}
\newcommand{\D}{\mathcal D}
\newcommand{\X}{\mathcal X}
\newcommand{\Rphys}{R_{\mathrm{phys}}}
\newcommand{\crit}{\epsilon_{\mathrm{crit},0}}
\newcommand{\Omk}{\Omega_K}
\newcommand{\LCDM}{\Lambda\mathrm{CDM}}
\newcommand{\Sthree}{S^3}
\newcommand{\Vol}{\operatorname{Vol}}
\newcommand{\diam}{\operatorname{diam}}
\newcommand{\Ob}{\Omega_{\rm b}}
\newcommand{\Oc}{\Omega_{\rm c}}
\newcommand{\Og}{\Omega_\gamma}
\newcommand{\Onu}{\Omega_\nu}
\newcommand{\OD}{\Omega_{\D}}
\newcommand{\OX}{\Omega_X}
\newcommand{\PX}{P_X}
\newcommand{\PD}{P_{\D}}
\newcommand{\Aacc}{\mathfrak A}
\newcommand{\Dacc}{\mathfrak D_X}
\newcommand{\Thetaexp}{\Theta}

\begin{document}

\title{Dark-sector residuals in compact \texorpdfstring{$S^3$}{S3} FLRW cosmology}

\author{Andrei~Preece\,\orcidlink{0009-0004-5978-795X}}
\email{andre.preece@web.de}
\affiliation{Independent Researcher, Berlin, Germany}

\author{Boris~Batenin\,\orcidlink{0009-0000-7027-8583}}
\email{batenin@web.de}
\affiliation{Independent Researcher, Berlin, Germany}

\date{May 2026}

\begin{abstract}
We formulate compact closed-FLRW cosmology on $\M=\RR\times S^3$ as a
conditional reconstruction problem for the gravitational dark residual.  The
assumptions are deliberately limited: general relativity, a closed homogeneous
and isotropic metric, an independently specified visible sector, and a specified
or reconstructed expansion history $E(a)=H(a)/H_0$.  Compact $S^3$ geometry
supplies finite volume, finite diameter, positive spatial curvature, causal
accessibility bounds, compact covariance kernels, and a discrete
scalar--vector--tensor harmonic bookkeeping system.  It does not by itself
imply dark matter, accelerated expansion, $w=-1$, or a cosmological constant.
The dark sector is instead defined by the covariant tensor residual
$T^\D_{\mu\nu}=c^4G_{\mu\nu}/(8\pi G)-T^{\rm vis}_{\mu\nu}$.  For a visible
baryon--photon--neutrino sector, the closed-FLRW background fixes
$\epsilon_\D(a)$ and $p_\D(a)$ uniquely from $E(a)$, $\Omega_K<0$, and the
visible densities.  Homogeneous distances therefore determine only the total
background residual; a pressureless clustering component and a smooth
accelerating component require perturbative observables and a specified closure
for perturbations.  We give the reconstruction theorem, invariant-closure
conditions, closed-distance formulae, compact $S^3$ spectral kernels, an
effective-Einstein interpretation for modified gravity, and a falsification
protocol using CMB, BAO, supernovae, cosmic chronometers, lensing, growth,
clusters, and galaxy dynamics.  Compact $S^3$ geometry supplies the finite
closed stage; the dark sector is the invariant residual target to be
reconstructed covariantly, closed invariantly, and tested observationally.
\end{abstract}

\maketitle

\section{Scope and central statement}

The model class is the compact closed-FLRW spacetime
\begin{equation}
\M=\RR\times S^3,\qquad
\dd s^2=-N^2(t)c^2\dd t^2+a^2(t)R_0^2\dd\Omega_3^2 ,
\label{eq:metric}
\end{equation}
with $a_0=1$ and $R_0$ the present curvature radius.  The dark residual is
defined by
\begin{equation}
T^{\D}_{\mu\nu}:=\frac{c^4}{8\pi G}G_{\mu\nu}-T^{\rm vis}_{\mu\nu}.
\label{eq:residual_tensor}
\end{equation}
This is a tensor definition inside general relativity.  It is not a gauge
choice, a particle model, a proof of a microscopic substance, or a topology-only
explanation.

\begin{nonclaim}[No topology-only dark sector]
The implications
\[
S^3\Longrightarrow \text{dark matter},\qquad
S^3\Longrightarrow \text{dark energy}
\]
are false.  Compactness fixes global spatial geometry and infrared spectral
structure.  It does not by itself imply clustering mass, $\ddot a>0$, $w=-1$,
or constant vacuum energy.
\end{nonclaim}

\begin{nonclaim}[No observationally large curvature assumption]
The assumption $\Sigma\simeq S^3$ does not require a curvature radius near the
present horizon.  The radius $R_0$ may be much larger than current cosmological
distance scales, and the reconstruction must be tested against the tight
CMB+BAO posterior on $\Omega_K$, which is consistent with flatness in standard
analyses \cite{Planck2018VI}.
\end{nonclaim}

\begin{nonclaim}[No tensor-vacuum mechanism]
No transverse-traceless vacuum source is used in the present construction.  A
future tensor source belongs in this framework only after an action or operator,
state, renormalization prescription, conserved stress tensor, perturbation
equations, and likelihood map are specified.
\end{nonclaim}

The strict implication proved here is
\begin{equation}
\begin{gathered}
S^3+\text{closed FLRW}+\text{GR}\\
+\text{visible sector}
+\text{specified or reconstructed }E(a)\\
\Longrightarrow
T^\D_{\mu\nu}\ \text{as a covariantly conserved residual}.
\end{gathered}
\label{eq:main_result}
\end{equation}
The physical interpretation of this residual is then tested by distances,
lensing, growth, cluster abundance, and galaxy dynamics.

The novelty is not the algebraic existence of Eq.~\eqref{eq:residual_tensor},
which follows once the Einstein equation is assumed.  The novelty is the
closed-geometry bookkeeping: the same compact background, visible-sector
closure, curvature convention, distance map, perturbation basis, and
falsification ledger are used to define the common target that particle dark
matter, dark energy, interacting sectors, or effective modified gravity must
reproduce.

\begin{remark}[Target residual]
The object reconstructed here is not a new substance and not a
topology-induced source.  It is the covariant target residual that any
dark-sector model, interacting sector, or effective modified-gravity model must
reproduce once the visible sector, closed geometry, and expansion history are
fixed.
\end{remark}

\section{Invariant kinematics and assumptions}

The metric signature is $(-,+,+,+)$.  The symbol $\epsilon$ denotes energy
density and $p$ denotes pressure, with the same dimensions.  Let $U^\mu$ be the
comoving four-velocity normalized by $U_\mu U^\mu=-c^2$.  In the coordinates of
Eq.~\eqref{eq:metric}, $U^\mu=(1/N,0,0,0)$.  The expansion scalar is defined
with this dimensional congruence,
\begin{equation}
\Thetaexp:=\nabla_\mu U^\mu ,
\label{eq:Theta}
\end{equation}
so that
\begin{equation}
H=\frac13\Thetaexp=\frac{\dot a}{Na}, \qquad E(a)=\frac{H(a)}{H_0}.
\label{eq:H_scalar}
\end{equation}
The present critical energy density and curvature parameter are
\begin{align}
\crit&=\frac{3c^2H_0^2}{8\pi G},&
\Rphys(a)&=aR_0,
\nonumber\\
\Omk&=-\frac{c^2}{H_0^2R_0^2}<0 .
\label{eq:ecrit_E}
\end{align}

The sufficient assumptions are:
\begin{enumerate}
\item $\Sigma\simeq S^3$ with the round constant-curvature metric in the
homogeneous background.
\item General relativity for the total stress tensor:
$G_{\mu\nu}=8\pi G T_{\mu\nu}/c^4$.
\item The visible stress tensor is specified independently and is covariantly
conserved unless an explicit interaction term is introduced.
\item The visible sector consists of baryons, photons, and neutrinos with
standard thermodynamic and kinetic closure.  Baryon density, neutrino masses,
and calibration parameters are sampled consistently in the same inference so that no
constraint is counted twice.
\item Distances, covariance kernels, lensing, growth, clusters, and galaxy
dynamics are evaluated in the same geometry and under the same priors.
\end{enumerate}
No claim of logical minimality is made.

\begin{proposition}[Invariant expansion function]
In the closed-FLRW class, $E(a)=H(a)/H_0$ is invariant under time
reparametrizations once the normalization $a_0=1$ is fixed.
\end{proposition}

\begin{proof}
The proper time of a comoving observer is $\dd\tau=N(t)\dd t$.  Hence
\[
H=\frac{1}{a}\frac{\dd a}{\dd\tau}
\]
is unchanged when $t$ is replaced by another time coordinate and the lapse is
transformed accordingly.  Division by the measured present value $H_0$ gives
$E(a)$.
\end{proof}

\section{Compact \texorpdfstring{$S^3$}{S3} geometry}

For the round three-sphere of radius $R$ \cite{Wald1984,HawkingEllis1973,doCarmo1992},
\begin{equation}
\Vol(S^3_R)=2\pi^2R^3,\qquad
{}^{(3)}\mathcal R=\frac{6}{R^2},\qquad
\diam(S^3_R)=\pi R .
\label{eq:s3_invariants}
\end{equation}
A geodesic ball of angular radius $\chi\in[0,\pi]$ has
\begin{equation}
V(\chi)=2\pi R^3\left(\chi-\frac12\sin2\chi\right),\qquad
A(\chi)=4\pi R^2\sin^2\chi .
\label{eq:ball}
\end{equation}

\begin{lemma}[Compact closure]
$S^3_R$ is compact, orientable, simply connected, boundaryless, and
geodesically complete.  Its maximum geodesic distance is $\pi R$.
\end{lemma}

\begin{proof}
The round three-sphere is the embedded set
$\{x\in\RR^4:|x|=R\}$ with the induced metric.  The distance is
$d(p,q)=R\arccos[(p\cdot q)/R^2]$.  Its maximum occurs at antipodal points.
Compact Riemannian manifolds are complete by Hopf--Rinow.
\end{proof}

For the ultrastatic metric, radial null propagation obeys $R\dd\chi=c\dd t$.
The accessible volume fraction after a time interval $T$ is therefore
\begin{equation}
\Gamma(T)=\frac{\theta-\frac12\sin2\theta}{\pi},\qquad
\theta=\min\left(\frac{cT}{R},\pi\right).
\label{eq:access}
\end{equation}
For the expanding metric, antipodal causal accessibility between $t_i$ and
$t_f$ is
\begin{equation}
\int_{t_i}^{t_f}\frac{c\,N(t)\dd t}{a(t)R_0}
=\int_{\tau_i}^{\tau_f}\frac{c\,\dd\tau}{a(\tau)R_0}\ge\pi .
\label{eq:conformal_access}
\end{equation}
This is causal reach, not a recurrence or topological echo.  In unquotiented
$S^3$ there are no nontrivial deck transformations.

\section{Compact spectral covariance and perturbative bookkeeping}
\label{sec:compact_spectral_covariance}

The compact spectral statement is independent of dark-sector microphysics.  On
a compact boundaryless Riemannian manifold, self-adjoint elliptic Laplace-type
operators have compact resolvent, real discrete spectrum, finite-dimensional
eigenspaces, and complete $L^2$ eigensections.  For the round three-sphere this
gives exact infrared discreteness while leaving the ultraviolet tower unbounded.

We label scalar harmonics by the principal $S^3$ level $n$,
\[
Y_{n\ell m},\qquad 0\leq \ell\leq n,\qquad -\ell\leq m\leq \ell.
\]
Here $n$ is the hyperspherical level and $\ell$ is an internal angular label; it
must not be confused with the observed sky multipole in a CMB map.  The scalar
Laplacian satisfies \cite{RubinOrdonez1984}
\begin{align}
-\nabla^2Y_{n\ell m}
&=\frac{n(n+2)}{R^2}Y_{n\ell m},
\nonumber\\
g_n^S&=(n+1)^2,\qquad n=0,1,2,\ldots .
\label{eq:scalar_spectrum}
\end{align}
For unit radius the isotropic scalar addition kernel is
\begin{equation}
\Xi_n(\gamma)=\frac{n+1}{2\pi^2}
\frac{\sin[(n+1)\gamma]}{\sin\gamma},\qquad
\gamma=d(x,x')/R .
\label{eq:addition_kernel}
\end{equation}
With the radius restored, the dimensionful kernel carries the factor $R^{-3}$.
An isotropic scalar covariance is therefore
\begin{equation}
C(\gamma)=\sum_{n=0}^{\infty} A_n\,\Xi_n(\gamma).
\label{eq:covariance}
\end{equation}
Compactness discretizes the infrared spectrum.  It does not impose an
ultraviolet cutoff, since $n$ is unbounded.

Linear perturbations on $\RR\times S^3$ require an $SO(4)$-adapted
scalar--vector--tensor basis rather than a continuum Fourier basis
\cite{RubinOrdonez1984,Higuchi1987,AbbottSchaefer1986,Jantzen1978,LindblomTaylorZhang2017}.
The eigenvalues below use the rough covariant Laplacian $-\nabla^a\nabla_a$ on
the round three-sphere.  A physical perturbation operator may contain additional
gauge, constraint, mass, or curvature-potential terms; such terms must be
specified by the model.

A scalar-derived longitudinal vector obeys, for $n\geq1$,
\begin{equation}
-\nabla^2(\nabla_aY_{n\ell m})
=\frac{n(n+2)-2}{R^2}\nabla_aY_{n\ell m}.
\label{eq:long_vector_eigen}
\end{equation}
The independent transverse vector harmonics satisfy
\begin{align}
\nabla^aY^{V,\sigma}_{n\ell m a}&=0,
\nonumber\\
-\nabla^2Y^{V,\sigma}_{n\ell m a}
&=\frac{n(n+2)-1}{R^2}Y^{V,\sigma}_{n\ell m a},
\qquad \sigma=1,2.
\label{eq:trans_vector_eigen}
\end{align}
with total transverse-vector degeneracy
\begin{equation}
g_n^{V_T}=2n(n+2),\qquad n\geq1.
\label{eq:vector_degeneracy}
\end{equation}
For $n=1$ these are the six Killing fields of the round three-sphere.

Symmetric rank-two perturbations decompose into scalar-derived, vector-derived,
and transverse-traceless parts.  The TT harmonics are defined by
\begin{equation}
\nabla^aY^{T,\sigma}_{n\ell m ab}=0,\qquad
g^{ab}Y^{T,\sigma}_{n\ell m ab}=0,\qquad \sigma=1,2.
\label{eq:tt_conditions}
\end{equation}
With the same rough-Laplacian convention,
\begin{equation}
-\nabla^2Y^{T,\sigma}_{n\ell m ab}
=\frac{n(n+2)-2}{R^2}Y^{T,\sigma}_{n\ell m ab},\qquad n\geq2,
\label{eq:tt_rough_eigen}
\end{equation}
and the total TT degeneracy is
\begin{equation}
g_n^{TT}=2(n-1)(n+3),\qquad n\geq2.
\label{eq:tt_degeneracy}
\end{equation}
Equivalently, with $j=n+1$, the rough eigenvalue is $(j^2-3)/R^2$ and the TT
degeneracy is $2(j^2-4)$ for $j\geq3$.

\begin{remark}[No scalar-spectrum substitution]
It is incorrect to assign the scalar eigenvalue $n(n+2)/R^2$ to all
perturbation sectors.  Vector and tensor harmonics carry curvature shifts, and
physical operators may carry further model-dependent terms.  A perturbative
dark-sector or modified-gravity model must therefore specify its actual
operator before claiming predictions for low-$n$ modes.
\end{remark}

The compact replacement rule for perturbative bookkeeping is
\begin{equation}
\int \frac{\dd^3 k}{(2\pi)^3}\,\mathcal F(k)
\quad\longrightarrow\quad
\frac{1}{\Vol(S^3_R)}
\sum_{n\in\mathcal N}\sum_{\alpha=1}^{g_n}\mathcal F_{n\alpha}.
\label{eq:compact_replacement_rule}
\end{equation}
The degeneracy index $\alpha$ belongs to the relevant scalar, vector, or tensor
sector, and $\mathcal N$ excludes homogeneous and gauge subspaces according to
the perturbation variable.  This is mode bookkeeping, not a claim that
compactness removes every infrared subtlety.  Gauge zero modes, state-dependent
long-wavelength effects, and de Sitter-like infrared questions remain
operator- and state-dependent.

For CMB temperature anisotropies, a scalar harmonic restricted to the
last-scattering sphere has the separated form
\begin{equation}
Y_{n\ell m}(\chi_*,\theta,\phi)=H_{n\ell}(\chi_*)Y_{\ell m}(\theta,\phi),
\label{eq:last_scattering_projection}
\end{equation}
so a closed-universe line-of-sight calculation replaces the flat Fourier
integral by a hyperspherical sum,
\begin{equation}
C_\ell^{TT}=\sum_{n\in\mathcal N_S,\,n\geq\ell}
\mathcal W_{n\ell}\,\mathcal P_{\mathcal R}(n)
\left|\Delta^T_{\ell n}\right|^2 .
\label{eq:closed_cmb_sum}
\end{equation}
This schematic expression is written up to the normalization convention used
for the hyperspherical harmonics, degeneracy weights, primordial spectrum, and
transfer functions.  Here $\ell$ is the observed sky multipole.  The transfer functions
$\Delta^T_{\ell n}$ must be computed in the same closed geometry and with the
same visible-sector and residual-sector parameters used in the background
likelihood; public Boltzmann codes such as CAMB or CLASS implement this when
$\Omega_K\neq0$ is specified \cite{Lewis2011,Lesgourgues2011}.
Equation~\eqref{eq:closed_cmb_sum} is deliberately conservative:
compactness supplies a discrete infrared spectrum, but the sign and size of any
low-$\ell$ change depend on the primordial spectrum, transfer functions, ISW
terms, reionization, foreground treatment, estimator, mask, and likelihood
\cite{Efstathiou2003,Planck2018VII}.

\section{Closed-FLRW background reconstruction}

With curvature kept on the geometric side, the Friedmann equation is
\cite{Wald1984,Peebles1993,DodelsonSchmidt2020,Weinberg2008}
\begin{equation}
H^2=\frac{8\pi G}{3c^2}\epsilon_{\rm tot}
-\frac{c^2}{a^2R_0^2}.
\label{eq:friedmann}
\end{equation}
Equivalently,
\begin{equation}
\frac{\epsilon_{\rm tot}(a)}{\crit}=E^2(a)-\Omk a^{-2}.
\label{eq:etot}
\end{equation}
Conservation gives
\begin{equation}
\dot\epsilon_{\rm tot}+3H(\epsilon_{\rm tot}+p_{\rm tot})=0,
\label{eq:conservation}
\end{equation}
and hence
\begin{equation}
\frac{p_{\rm tot}(a)}{\crit}
=-\frac13\left[\frac{\dd E^2}{\dd\ln a}+3E^2(a)-\Omk a^{-2}\right].
\label{eq:ptot}
\end{equation}
The deceleration scalar is
\begin{equation}
q(a):=-\frac{\ddot a\,a}{\dot a^2}
=-1-\frac12\frac{\dd\ln E^2}{\dd\ln a}.
\label{eq:q}
\end{equation}
The acceleration equation is
\begin{equation}
\frac{\ddot a}{a}=-\frac{4\pi G}{3c^2}
\left(\epsilon_{\rm tot}+3p_{\rm tot}\right).
\label{eq:accel}
\end{equation}
Thus accelerated expansion is the scalar inequality
\begin{equation}
\epsilon_{\rm tot}+3p_{\rm tot}<0.
\label{eq:accel_condition}
\end{equation}
Define
\begin{equation}
\Aacc(a):=-\frac{1}{2}\frac{\epsilon_{\rm tot}+3p_{\rm tot}}{\crit}
=-q(a)E^2(a).
\label{eq:Aacc}
\end{equation}
Acceleration is equivalent to $\Aacc(a)>0$.

\section{Visible-sector closure}

The visible tensor is
\begin{equation}
T^{\rm vis}_{\mu\nu}=T^b_{\mu\nu}+T^\gamma_{\mu\nu}+T^\nu_{\mu\nu}.
\label{eq:visible}
\end{equation}
On the homogeneous background,
\begin{align}
\frac{\epsilon_{\rm vis}(a)}{\crit}
&=\Ob a^{-3}+\Og a^{-4}+\Onu(a),
\nonumber\\
\frac{p_{\rm vis}(a)}{\crit}
&=\frac13\Og a^{-4}+P_\nu(a).
\label{eq:visible_background}
\end{align}

The thermodynamic photon and neutrino densities are evaluated in the standard
local kinetic description used by cosmological Boltzmann solvers.  Exact
finite-volume mode discreteness of quantum fields on a compact $S^3$ would be a
separate state- and operator-dependent correction to the visible stress tensor;
it is not part of the dark residual unless it is explicitly added to
$T^{\rm vis}_{\mu\nu}$.

The baryon term is the rest-energy density of conserved nonrelativistic baryon
number,
\begin{equation}
\epsilon_b(a)=\crit\,\Ob a^{-3}.
\label{eq:baryon_closure}
\end{equation}
The photon term is fixed by the observed CMB temperature \cite{Fixsen2009},
\begin{align}
\epsilon_\gamma(a)
&=\frac{\pi^2}{15}
\frac{(k_{\rm B}T_{\gamma0})^4}{\hbar^3c^3}\,a^{-4}
=\crit\,\Og a^{-4},
\nonumber\\
p_\gamma(a)&=\frac13\epsilon_\gamma(a).
\label{eq:photon_closure}
\end{align}

\subsection{Massive-neutrino functions}

The neutrino temperature and phase-space form follow from standard weak
decoupling, entropy conservation, and collisionless propagation
\cite{DodelsonSchmidt2020,Mukhanov2005,LesgourguesPastor2006,LesgourguesPastor2012}.
The absolute neutrino mass scale is not fixed by thermodynamics alone; it is an
input or nuisance parameter constrained by laboratory and cosmological priors.

At temperatures below the $W$-boson mass, the interaction rate scales as
\begin{equation}
\Gamma_\nu(T)\simeq C_\nu G_F^2T^5,\qquad C_\nu=\mathcal O(1),
\label{eq:nu_rate}
\end{equation}
while during radiation domination
\begin{equation}
H(T)=1.66\,g_*^{1/2}\frac{T^2}{M_{\rm Pl}}.
\label{eq:radiation_H}
\end{equation}
Freeze-out occurs when $\Gamma_\nu\simeq H$, giving
\begin{equation}
T_{\rm dec}\simeq
\left(\frac{1.66\,g_*^{1/2}}{C_\nu G_F^2M_{\rm Pl}}\right)^{1/3}
=\mathcal O(1)\ {\rm MeV}
\label{eq:Tdec}
\end{equation}
for $g_*=10.75$ near the MeV epoch.  Neutrinos therefore decouple while
ultrarelativistic.

After decoupling, entropy transfer from $e^\pm$ annihilation heats photons
relative to neutrinos.  In the instantaneous-decoupling approximation,
\begin{equation}
g_{*s,{\rm before}}^{\rm em}=2+\frac{7}{8}\times 4=\frac{11}{2},\qquad
g_{*s,{\rm after}}^{\rm em}=2,
\end{equation}
and
\begin{equation}
T_{\nu0}=\left(\frac{4}{11}\right)^{1/3}T_{\gamma0}\simeq1.95\ {\rm K}.
\label{eq:Tnu0}
\end{equation}
Non-instantaneous decoupling, finite-temperature QED corrections, and flavor
oscillations are commonly summarized by the Standard-Model value
$N_{\rm eff}\simeq3.044$ \cite{Mangano2005,deSalasPastor2016,Froustey2020}.
The analytic formulae below use the instantaneous-decoupling Fermi--Dirac form.
In numerical likelihoods this ideal distribution should be replaced by
Boltzmann-solver neutrino transfer functions or tabulated distorted spectra
consistent with the adopted $N_{\rm eff}$.  The tensor residual formulae below
are unchanged.

In Eqs.~\eqref{eq:eps_nui}--\eqref{eq:p_nui} we use natural units
$\hbar=c=k_{\rm B}=1$.  For each mass eigenstate $m_i$, let
\begin{equation}
f_0(q)=\frac{1}{\exp(q/T_{\nu0})+1}
\label{eq:f0_nu}
\end{equation}
be the frozen occupation number as a function of comoving momentum $q=ap$,
with vanishing chemical potential.  With degeneracy $g_i=2$ for the neutrino
plus antineutrino helicity states of one mass eigenstate,
\begin{align}
\epsilon_{\nu i}(a)&=
\frac{g_i}{2\pi^2a^4}
\int_0^\infty \dd q\,q^2
\sqrt{q^2+a^2m_i^2}\, f_0(q),
\label{eq:eps_nui}\\
p_{\nu i}(a)&=
\frac{g_i}{6\pi^2a^4}
\int_0^\infty \dd q\,
\frac{q^4}{\sqrt{q^2+a^2m_i^2}}\, f_0(q).
\label{eq:p_nui}
\end{align}
The total neutrino contribution is
\begin{equation}
\epsilon_\nu(a)=\sum_{i=1}^{3}\epsilon_{\nu i}(a),\qquad
p_\nu(a)=\sum_{i=1}^{3}p_{\nu i}(a),
\end{equation}
and
\begin{equation}
\Omega_\nu(a)=\frac{\epsilon_\nu(a)}{\crit},\qquad
P_\nu(a)=\frac{p_\nu(a)}{\crit}.
\label{eq:Onu_Pnu}
\end{equation}
The limiting checks are
\begin{align}
a m_i\ll q:\quad&
\epsilon_{\nu i}\propto a^{-4},\qquad
p_{\nu i}=\frac13\epsilon_{\nu i},
\label{eq:nu_rel_limit}\\
a m_i\gg q:\quad&
\epsilon_{\nu i}\simeq m_i n_{\nu i}\propto a^{-3},\qquad
p_{\nu i}\ll\epsilon_{\nu i}.
\label{eq:nu_nonrel_limit}
\end{align}
At late times, for nonrelativistic neutrinos,
\begin{equation}
\Omega_{\nu0}h^2\simeq
\frac{\sum_i m_i}{93.14\ {\rm eV}} .
\label{eq:Onu_mass_sum}
\end{equation}
Since the post-decoupling distribution is collisionless,
\begin{equation}
\dot\epsilon_\nu+3H(\epsilon_\nu+p_\nu)=0.
\label{eq:nu_conservation}
\end{equation}

\section{GR dark residual}

\begin{theorem}[GR residual theorem]
Within general relativity, assume the total stress tensor obeys the Einstein
equation and the visible stress tensor is covariantly conserved,
$\nabla_\mu T^{{\rm vis}\,\mu\nu}=0$.  Then the tensor
$T^\D_{\mu\nu}$ defined by Eq.~\eqref{eq:residual_tensor} is symmetric and
covariantly conserved.  On the closed-FLRW background it has the perfect-fluid
form
\[
T^{\D}_{\hat\mu\hat\nu}
=\operatorname{diag}(\epsilon_{\D},p_{\D},p_{\D},p_{\D})
\]
in the comoving orthonormal frame, with
\begin{align}
\frac{\epsilon_{\D}(a)}{\crit}
&=E^2(a)-\Omk a^{-2}
-\Ob a^{-3}-\Og a^{-4}-\Onu(a),
\label{eq:dark_energy_residual}\\
\frac{p_{\D}(a)}{\crit}
&=-\frac13\left[\frac{\dd E^2}{\dd\ln a}
+3E^2(a)-\Omk a^{-2}\right]
\nonumber\\
&\quad -\frac13\Og a^{-4}-P_\nu(a).
\label{eq:dark_pressure_residual}
\end{align}
\end{theorem}

\begin{proof}
Both $G_{\mu\nu}$ and $T^{\rm vis}_{\mu\nu}$ are symmetric rank-two tensors;
their linear combination is therefore a symmetric rank-two tensor.  The
twice-contracted Bianchi identity gives $\nabla_\mu G^{\mu\nu}=0$, hence
\[
\nabla_\mu T^{\D\,\mu\nu}
=\frac{c^4}{8\pi G}\nabla_\mu G^{\mu\nu}
-\nabla_\mu T^{{\rm vis}\,\mu\nu}=0 .
\]
Homogeneity and isotropy force the background part of any covariantly defined
symmetric tensor to be a linear combination of $U_\mu U_\nu$ and the spatial
projector $h_{\mu\nu}:=g_{\mu\nu}+c^{-2}U_\mu U_\nu$, by Schur's lemma applied
to the isotropy group $\mathrm{SO}(3)$; it is therefore a perfect fluid.  The
$00$ Einstein equation gives Eq.~\eqref{eq:etot}.  Subtracting
Eq.~\eqref{eq:visible_background} gives Eq.~\eqref{eq:dark_energy_residual}.
For pressure, Eq.~\eqref{eq:conservation} gives
$p_{\rm tot}=-\epsilon_{\rm tot}
-\frac13\dd\epsilon_{\rm tot}/\dd\ln a$, which is Eq.~\eqref{eq:ptot}.
Subtracting the visible pressure gives Eq.~\eqref{eq:dark_pressure_residual}.
\end{proof}

\begin{remark}[Curvature is not relabeled as a dark fluid]
The curvature term is kept as geometry through $\Omega_K$ and is not counted as
dark matter or dark energy.  A curvature-only explanation would require the
residual reconstructed by Eqs.~\eqref{eq:dark_energy_residual} and
\eqref{eq:dark_pressure_residual} to vanish or become observationally
unnecessary in a global likelihood analysis.
\end{remark}

The residual equation of state, where defined, is
\begin{equation}
w_{\D}(a)=\frac{p_{\D}(a)}{\epsilon_{\D}(a)}
=\frac{\PD(a)}{\OD(a)} .
\label{eq:wD}
\end{equation}

\begin{remark}[Residual is not a microphysical explanation]
Equation~\eqref{eq:residual_tensor} reconstructs what is missing from the
specified visible sector within GR.  It does not distinguish particle dark
matter, vacuum energy, modified gravity written as an effective tensor, or an
interacting dark sector.  That distinction requires perturbations and model
comparison.
\end{remark}

\section{Clustering residual and smooth residual}

The homogeneous background cannot uniquely split the residual into dark matter
and dark energy.

\begin{theorem}[Background degeneracy]
On an FLRW background, any residual density $\epsilon_\D(a)$ can be decomposed
as $\epsilon_\D(a)=\epsilon_c(a)+\epsilon_X(a)$ in infinitely many ways unless
additional information fixes clustering, equation of state, sound speed,
anisotropic stress, or interactions.  Homogeneous distances alone do not
identify dark matter.
\end{theorem}

\begin{proof}
Choose any differentiable function $\epsilon_c(a)$ and define
$\epsilon_X=\epsilon_\D-\epsilon_c$.  The same total homogeneous residual is
obtained.  Distances see only the total background expansion.  Perturbative
observables are required to separate components.
\end{proof}

The phenomenological split used for comparison with standard cosmological
likelihoods is
\begin{align}
\Omega_c(a)&=\Omega_{c0}a^{-3},\qquad p_c=0,
\label{eq:cold_dm}\\
\Omega_X(a)&=E^2(a)-\Omk a^{-2}
-\Ob a^{-3}-\Og a^{-4}-\Onu(a)
\nonumber\\
&\quad -\Omega_{c0}a^{-3},
\label{eq:OmegaX}\\
P_X(a)&=-\frac13\left[\frac{\dd E^2}{\dd\ln a}
+3E^2(a)-\Omk a^{-2}\right]
\nonumber\\
&\quad -\frac13\Og a^{-4}-P_\nu(a),
\label{eq:PX}\\
w_X(a)&=\frac{P_X(a)}{\Omega_X(a)}.
\label{eq:wX}
\end{align}
Here $\Omega_{c0}$ is not determined by the background alone.  Within GR and a
specified perturbation closure it is constrained by CMB peak structure, CMB
lensing, large-scale-structure growth, weak lensing, cluster abundance, and
galaxy-scale dynamical mass diagnostics
\cite{Planck2018VI,ACTDR62024,DESY32022,BulletCluster2006,SPARC2016}.

For the smooth residual define
\begin{equation}
\Dacc(a):=\Omega_X(a)+3P_X(a).
\label{eq:DX}
\end{equation}
The smooth component carries a negative contribution to the total
$\epsilon_{\rm tot}+3p_{\rm tot}$ when $\Dacc(a)<0$; the sign of
the total acceleration is fixed only after matter and radiation are included in
Eq.~\eqref{eq:accel}.

A cosmological constant is the special case
\begin{align}
T^X_{\mu\nu}&=-\epsilon_\Lambda g_{\mu\nu},
&
\Omega_X(a)&=\Omega_\Lambda,
\nonumber\\
P_X&=-\Omega_\Lambda,
&
w_X&=-1,
&
\Dacc&=-2\Omega_\Lambda .
\label{eq:Lambda_case}
\end{align}
A dynamical dark-energy reconstruction requires a statistically significant
failure of Eq.~\eqref{eq:Lambda_case} after curvature, neutrino masses,
calibration, and selection effects are marginalized
\cite{Linder2003,DESIDR2DE2025}.  Current DESI DR2 combinations motivate such
tests, but they do not by themselves establish a physical departure from
$\Lambda$CDM \cite{DESIDR2BAO2025,DESIDR2DE2025}.

\section{Observable distances and perturbative invariants}

For $z=a^{-1}-1$,
\begin{equation}
D_C(z)=c\int_0^z\frac{\dd z'}{H(z')},\qquad
\chi(z)=\frac{D_C(z)}{R_0}.
\label{eq:DC_chi}
\end{equation}
The closed transverse comoving distance and the standard angular and luminosity
distances are
\begin{align}
D_M(z)&=R_0\sin\chi(z),
\nonumber\\
D_A(z)&=\frac{D_M(z)}{1+z},
&
D_L(z)&=(1+z)D_M(z).
\label{eq:closed_distances}
\end{align}
BAO, supernovae, and CMB distances constrain $E(a)$ and $\Omega_K$ through
these scalar distances
\cite{Planck2018VI,DESIDR2BAO2025,DESIDR2Lya2025,
PantheonPlus2022,PantheonPlusData2022}.  Common BAO
combinations are
\begin{align}
D_H(z)&=\frac{c}{H(z)},
&
D_V(z)&=\left[zD_M^2(z)D_H(z)\right]^{1/3},
\nonumber\\
F_{\rm AP}(z)&=\frac{D_M(z)H(z)}{c}.
\label{eq:bao_invariants}
\end{align}
Etherington reciprocity gives
\begin{equation}
D_L(z)=(1+z)^2D_A(z)
\label{eq:etherington}
\end{equation}
when photons travel on null geodesics and photon number is conserved
\cite{Etherington1933}.

The linear growth observables are
\begin{equation}
f(a)=\frac{\dd\ln D_+(a)}{\dd\ln a},\qquad f\sigma_8(z),
\label{eq:growth}
\end{equation}
as measured by redshift-space distortions.  Weak lensing and CMB lensing
measure projected Weyl-potential combinations along null geodesics
\cite{ACTDR62024,DESY32022}.  Merging clusters supply a complementary
diagnostic: offsets between lensing mass and the dominant X-ray plasma are
evidence for a gravitating component displaced from the main baryonic gas, not
by themselves a proof of a particular particle model \cite{BulletCluster2006}.


\section{Complete invariant closure of the dark-sector target}
\label{sec:complete_invariant_closure}

The phrase ``closing the dark sector'' has two logically different meanings.
The first is an exact reconstruction of the stress tensor required by the metric
and visible sector.  The second is a microphysical explanation of that tensor in
terms of particles, fields, vacuum energy, interactions, or modified gravity.
Only the first is model-independent.  The second requires extra dynamical
structure and cannot be obtained from topology or background geometry alone.

\begin{definition}[Invariant residual closure]
Fix the closed-FLRW background, the visible stress tensor, and a perturbative
order.  The dark residual is invariantly closed at that order if, for every
scalar, vector, and tensor harmonic coefficient retained in the calculation,
\begin{equation}
T^{\D}_{(r)\mu\nu}
=\frac{c^4}{8\pi G}G_{(r)\mu\nu}-T^{{\rm vis}}_{(r)\mu\nu},
\qquad r=0,1,\ldots,N ,
\label{eq:order_by_order_residual}
\end{equation}
is specified in a gauge-invariant representation, and the same variables are
used in the distance, lensing, growth, CMB, cluster, and galaxy likelihoods.
Here $r=0$ denotes the homogeneous background and $r\geq1$ denotes perturbative
orders.
\end{definition}

At linear scalar order one may choose the Bardeen potentials, matter velocity
potentials, density contrasts, and anisotropic stresses.  In the closed model
each of these is expanded in hyperspherical harmonics rather than Fourier
modes.  Let $\mathcal T^\D_{n\ell m}$ denote the scalar coefficient tuple
$\{\delta\epsilon^\D_{n\ell m},\delta p^\D_{n\ell m},
(\epsilon_\D+p_\D)V^\D_{n\ell m},\Pi^\D_{n\ell m}\}$, and let
$\mathcal G_{n\ell m}$ and $\mathcal T^{\rm vis}_{n\ell m}$ denote the
corresponding Einstein and visible-sector tuples.  Then
\begin{equation}
\mathcal T^\D_{n\ell m}
=\frac{c^4}{8\pi G}\mathcal G_{n\ell m}
-\mathcal T^{\rm vis}_{n\ell m}.
\label{eq:scalar_residual_coefficients}
\end{equation}
Equation~\eqref{eq:scalar_residual_coefficients} is not a new field equation.
It is the harmonic version of the tensor identity
Eq.~\eqref{eq:residual_tensor}.  Vector and tensor coefficients are obtained in
the corresponding transverse-vector and transverse-traceless bases of
Sec.~\ref{sec:compact_spectral_covariance}.

\begin{theorem}[No purely geometric closure of the dark sector]
No collection of invariants constructed only from the topology of $S^3$, the
round spatial curvature radius, the compact Laplace spectra, the homogeneous
distance map, or their algebraic combinations can determine a unique physical
split of $T^\D_{\mu\nu}$ into dark matter and dark energy.
\end{theorem}

\begin{proof}
Those invariants determine the stage on which the metric and perturbations are
expanded.  They do not determine an equation of state, sound speed, anisotropic
stress, interaction current, rest frame, or clustering law for an unknown
component.  On the homogeneous background the degeneracy theorem already gives
infinitely many decompositions
$\epsilon_\D=\epsilon_c+\epsilon_X$.  At perturbative order the same residual
metric response can be represented by different combinations of density
contrast, pressure perturbation, anisotropic stress, and modified-gravity
response functions unless a dynamical closure is imposed.  Therefore geometry
alone cannot supply the missing microphysics.
\end{proof}

\begin{theorem}[Complete covariant target closure]
Within GR, given the metric to the chosen perturbative order and a separately
specified visible sector, Eq.~\eqref{eq:order_by_order_residual} uniquely
closes the dark residual as a tensor.  This is the maximal model-independent
closure: every particle dark-sector model, interacting-sector model, vacuum
model, or effective modified-gravity model must reproduce these residual
coefficients in the same closed geometry.
\end{theorem}

\begin{proof}
The Einstein tensor is determined by the metric order by order.  The visible
stress tensor is specified independently.  Their difference is a unique tensor
coefficient in each scalar, vector, and tensor harmonic sector.  Gauge-invariant
variables give the same tensor coefficients in any coordinate representation.
Any model producing the same metric and visible sector must therefore produce
the same residual tensor, or the same effective tensor if it is written in
Einstein form.
\end{proof}

A compact way to display the complete invariant target is
\begin{equation}
\mathcal C_{\rm inv}=
\left[
\Omega_K,\ E(a),\ T^{\rm vis}_{\mu\nu};
\ T^\D_{\mu\nu};
\ \{\delta T^\D_{n\alpha}\}_{S,V,T};
\ \mathcal R_{\rm dyn}
\right].
\label{eq:complete_invariant_target}
\end{equation}
The first entries fix the closed background and visible sector.  The middle
entries are the exact background and perturbative residuals.  The final entry
$\mathcal R_{\rm dyn}$ is not determined by geometry; it denotes the dynamical
response required to interpret the residual as dust, a smooth component,
vacuum energy, an interacting sector, or modified gravity.  In a likelihood
analysis the corresponding posterior has the schematic invariant factorization
\begin{multline}
\mathcal P(\mathcal C_{\rm inv}\mid\mathcal D)
\propto
\pi(\mathcal C_{\rm inv})\,
\mathcal L_{\rm CMB}\mathcal L_{\rm BAO}\mathcal L_{\rm SN}\\
{}\times
\mathcal L_{H(z)}\mathcal L_{\rm lens}
\mathcal L_{\rm growth}\mathcal L_{\rm cl}\mathcal L_{\rm gal}.
\label{eq:closure_posterior}
\end{multline}
where all likelihoods use the same curvature, visible-sector priors, and
closed-harmonic perturbation convention.  Observational closure means that the
same posterior samples pass all of these tests without moving curvature,
neutrino mass, calibration, or selection effects between inconsistent priors.

The invariant closure filters that have physical content are the following.

\begin{enumerate}
\item \textit{Pressureless clustering filter.}  A component $c$ qualifies as
cold clustering matter only if
\begin{align}
T^c_{\mu\nu}
&=\frac{\epsilon_c}{c^2}u^c_\mu u^c_\nu,
&
u^c_\mu u_c^\mu&=-c^2,
\nonumber\\
p_c&=0,\qquad
\Pi^c_{\mu\nu}=0,
\nonumber\\
\nabla_\mu T^{c\,\mu\nu}&=Q_c^\nu .
\label{eq:dust_filter}
\end{align}
with a stated interaction current $Q_c^\nu$ and with $Q_c^\nu=0$ in the
separately conserved case.  In the noninteracting limit the flow is geodesic
where $\epsilon_c\ne0$.  This filter is tested by CMB peak structure, lensing,
growth, cluster abundance, and galaxy dynamics.

\item \textit{Smooth accelerating filter.}  The residual left after the
clustering component is
$T^X_{\mu\nu}=T^\D_{\mu\nu}-T^c_{\mu\nu}$.  It contributes negatively to the
Raychaudhuri source when
\begin{equation}
\Dacc(a)=\Omega_X(a)+3P_X(a)<0 .
\label{eq:smooth_acceleration_filter}
\end{equation}
Calling it smooth is an observational and dynamical statement: the perturbation
closure must specify its rest-frame sound speed, anisotropic stress, and any
interaction current, or show that its clustering is negligible in the relevant
likelihood.

\item \textit{Cosmological-constant filter.}  The smooth component is exactly a
cosmological constant if and only if
\begin{equation}
T^X_{\mu\nu}=-\epsilon_\Lambda g_{\mu\nu},\qquad
\nabla_\alpha\epsilon_\Lambda=0,
\label{eq:lambda_filter}
\end{equation}
so that all gauge-invariant perturbations of this component vanish.  In the
background language this is the joint condition
$\Omega_X=\Omega_\Lambda$, $P_X=-\Omega_\Lambda$, and $w_X=-1$ at all $a$.

\item \textit{Modified-gravity filter.}  If the residual is moved to the
right-hand side of an effective Einstein equation, the background conditions
above are necessary but insufficient.  The model must also provide invariant
metric-response functions, or their closed-harmonic equivalents, that reproduce
both the dynamical potential and the Weyl potential probed by growth and
lensing.  In phenomenological notation these are often represented by functions
such as $\mu(a,n)$, $\Sigma(a,n)$, and a gravitational-slip variable.
\end{enumerate}

Thus the exhaustive invariant answer is narrow but precise.  Topology and
compact spectra close the geometry.  The Einstein residual closes the required
stress tensor.  A physical dark-sector explanation is closed only after one of
the dynamical filters above supplies the missing perturbative response and
passes the joint likelihood.

\section{Diagnostic ledger}

The following quantities form a compact diagnostic ledger.  They are not
independent axioms; they are invariant or operational quantities used to
reconstruct and test the residual.

\begin{enumerate}
\item $\Sigma\simeq S^3$: compact closed spatial topology.
\item $\partial S^3=\varnothing$: no spatial boundary source in the unquotiented model.
\item $\pi_1(S^3)=0$: no holonomy images or deck transformations.
\item $V_3=2\pi^2R^3$: finite spatial volume at fixed scale factor.
\item $\diam(S^3_R)=\pi R$: finite diameter.
\item ${}^{(3)}\mathcal R=6/R^2$: positive spatial curvature scalar.
\item $\Omk=-c^2/(H_0^2R_0^2)$: closed curvature parameter.
\item $\Thetaexp=\nabla_\mu U^\mu$: invariant expansion scalar.
\item $E(a)=H/H_0$: dimensionless expansion history.
\item $q(a)$ and $\Aacc(a)$: acceleration diagnostics.
\item $\epsilon_{\rm tot}/\crit=E^2-\Omk a^{-2}$: total density reconstruction.
\item $p_{\rm tot}/\crit$: total pressure reconstruction.
\item $T^\D_{\mu\nu}$: covariant dark residual.
\item $\epsilon_\D(a),p_\D(a)$: homogeneous residual scalars.
\item $\nabla_\mu T^{\D\,\mu\nu}=0$: Bianchi consistency under visible conservation.
\item $\Omega_bh^2$: baryon density constrained by CMB and BBN.
\item $T_{\gamma0}$: photon temperature fixing $\Omega_\gamma$.
\item $T_{\nu0}$ and $f_0(q)$: neutrino thermal closure.
\item $\sum_i m_i$: neutrino mass-sum control parameter.
\item $\Omega_c a^{-3},p_c=0$: pressureless clustering ansatz.
\item CMB peak structure: matter loading and equality diagnostics.
\item $C_L^{\phi\phi}$: CMB lensing mass map.
\item Weak shear $\xi_\pm$: projected late-time mass correlations.
\item $f\sigma_8(z)$: growth-rate observable.
\item Cluster abundance: mass-function test.
\item Galaxy dynamics: small-scale mass-discrepancy diagnostic.
\item $\Omega_X(a),P_X(a)$: smooth residual after clustering matter.
\item $w_X(a)$: smooth equation-of-state reconstruction.
\item $\Dacc=\Omega_X+3P_X$: smooth acceleration factor.
\item $D_M,D_A,D_L,D_H,D_V,F_{\rm AP}$: closed-distance and BAO observables.
\item $\theta_*=r_s(z_*)/D_M(z_*)$: CMB acoustic-angle constraint.
\item $\lambda_n=n(n+2)/R^2$: compact scalar spectrum.
\item $g_n^S=(n+1)^2$: scalar mode count.
\item $\Xi_n(\gamma)$: addition-kernel invariant.
\item $C(\gamma)=\sum_n A_n\Xi_n$: compact scalar covariance.
\item $g_n^{V_T}=2n(n+2)$: transverse-vector mode count.
\item $g_n^{TT}=2(n-1)(n+3)$: transverse-traceless tensor mode count.
\item Eq.~\eqref{eq:compact_replacement_rule}: compact perturbation replacement rule.
\item Eq.~\eqref{eq:closed_cmb_sum}: closed-geometry CMB transfer-sum interface.
\item Eq.~\eqref{eq:complete_invariant_target}: minimal complete invariant target.
\item Eq.~\eqref{eq:order_by_order_residual}: order-by-order invariant residual closure.
\item Eq.~\eqref{eq:closure_posterior}: joint posterior closure criterion.
\item Eqs.~\eqref{eq:dust_filter}--\eqref{eq:lambda_filter}: clustering, acceleration, and cosmological-constant filters.
\item Posterior predictive residuals: falsification beyond best fit.
\item Flat metric limit $R_0\to\infty$: local reduction to standard FLRW distances.
\end{enumerate}

\section{Data-facing reconstruction protocol}

\paragraph*{Invariant reconstruction protocol.}
\begin{enumerate}
\item Choose the geometry, field equation, priors, and $E(a)$ basis.
\item Reconstruct $E(a)$ and $\Omega_K$ from CMB, BAO, supernovae, and
chronometers \cite{Moresco2022}.
\item Fix $\Omega_bh^2$, $\Omega_\gamma$, and the neutrino functions
$\Omega_\nu(a),P_\nu(a)$ conditional on the chosen mass spectrum, BBN/CMB
constraints, laboratory priors, and jointly sampled cosmological nuisance priors
\cite{Cooke2018,SchrammTurner1998}.
\item Compute $\epsilon_\D(a)$ and $p_\D(a)$ from
Eqs.~\eqref{eq:dark_energy_residual} and \eqref{eq:dark_pressure_residual}.
\item Apply the invariant closure filters of
Sec.~\ref{sec:complete_invariant_closure}: residual-only closure, dust-like
clustering closure, smooth-acceleration closure, cosmological-constant closure,
or effective modified-gravity closure.
\item Infer $\Omega_{c0}$ from perturbative observables: CMB peaks, CMB lensing,
weak lensing, $f\sigma_8$, clusters, and galaxy dynamics.
\item Compute $\Omega_X(a)$, $P_X(a)$, $w_X(a)$, and $\Dacc(a)$ from
Eqs.~\eqref{eq:OmegaX}--\eqref{eq:DX}.
\item Compare $\Omega_X=\mathrm{const}$, $w_X=-1$, evolving $w(z)$, interacting
dark sectors, modified-gravity effective tensors, and curvature-only
explanations under the same likelihoods.
\item Build compact $S^3$ scalar, vector, and tensor mode expansions using
Eq.~\eqref{eq:compact_replacement_rule}; compare low-multipole CMB covariance
with the same curvature posterior.
\item Report posterior predictive checks, evidence or information criteria, and
the local flat-metric limit $R_0\to\infty$.
\end{enumerate}

Finite noisy data reconstruct a posterior over $E(a)$, not a unique curve.  The
uniqueness statements are conditional on a specified geometry, visible sector,
and exact or posterior-sampled expansion history.

\section{Formal consequences}

\begin{theorem}[Invariant reconstruction theorem]
Given a closed-FLRW geometry, the GR field equation, a covariantly conserved
visible stress tensor, and an expansion history $E(a)$ with curvature
$\Omega_K$, the residual $T^\D_{\mu\nu}$ is uniquely determined as a tensor by
Eq.~\eqref{eq:residual_tensor}.  Its background energy density and pressure are
the scalars in Eqs.~\eqref{eq:dark_energy_residual} and
\eqref{eq:dark_pressure_residual}.
\end{theorem}

\begin{proof}
The metric determines the Einstein tensor.  The visible tensor is specified
independently.  Their difference is a unique tensor.  FLRW symmetry reduces its
background part to two scalar functions, computed above.
\end{proof}

\begin{theorem}[Perturbative separation theorem]
Within GR and a specified perturbation closure, a pressureless clustering
component and a smooth accelerating component cannot be distinguished by
homogeneous distances alone.  They can be separated only by perturbative
observables such as CMB peak structure, lensing, growth, cluster abundance, and
galaxy-scale dynamical mass diagnostics.
\end{theorem}

\begin{proof}
Homogeneous distances constrain $E(a)$ and therefore only the total homogeneous
density and pressure.  Perturbations depend on clustering, sound speed,
anisotropic stress, and gravitational potentials.  These quantities affect
lensing and growth differently from a smooth cosmological-constant-like
component.  Hence the split is observationally meaningful only after
perturbative data and a perturbation model are specified.
\end{proof}

\section{Effective interpretation in modified gravity}

The residual construction is not tied to a particle interpretation.  It also
gives the target effective tensor for any metric modified-gravity model that can
be written in Einstein form with separately conserved visible matter.

\begin{proposition}[Effective Einstein form]
Suppose a diffeomorphism-covariant metric theory has field equations that can be
rearranged as
\begin{equation}
G_{\mu\nu}
=\frac{8\pi G}{c^4}\left(T^{\rm vis}_{\mu\nu}
+T^{\rm eff}_{\mu\nu}\right),
\label{eq:effective_einstein_form}
\end{equation}
with $\nabla_\mu T^{{\rm vis}\,\mu\nu}=0$.  Then the effective stress tensor
required by the modified theory is exactly the covariant residual,
\begin{equation}
T^{\rm eff}_{\mu\nu}=T^\D_{\mu\nu}.
\label{eq:Teff_equals_TD}
\end{equation}
Consequently its homogeneous energy density and pressure must equal
Eqs.~\eqref{eq:dark_energy_residual} and \eqref{eq:dark_pressure_residual}.
\end{proposition}

\begin{proof}
Equation~\eqref{eq:effective_einstein_form} gives
\[
T^{\rm eff}_{\mu\nu}
=\frac{c^4}{8\pi G}G_{\mu\nu}-T^{\rm vis}_{\mu\nu},
\]
which is precisely Eq.~\eqref{eq:residual_tensor}.  Conservation follows from
the Bianchi identity and visible-sector conservation exactly as in the GR
residual theorem.
\end{proof}

The proposition is an identity, not a new dynamical model.  For example, an
$f(R)$ model can be moved into effective Einstein form by placing the extra
curvature terms on the right-hand side \cite{HuSawicki2007}.  The homogeneous
target is then fixed:
\[
\epsilon_{\rm eff}(a)=\epsilon_\D(a),\qquad
p_{\rm eff}(a)=p_\D(a).
\]
A model that cannot reproduce these functions for the reconstructed $E(a)$ and
$\Omega_K$ fails at the homogeneous level.

The equivalence is conditional.  If visible matter is nonminimally coupled so
that $\nabla_\mu T^{{\rm vis}\,\mu\nu}\ne0$, or if photons do not follow the
metric null geodesics used in Eq.~\eqref{eq:closed_distances}, then the present
reconstruction theorem must be modified.  The interaction current, altered
distance law, and perturbation equations must then be stated explicitly.

At the perturbative level, matching the background residual is insufficient.  A
viable modified-gravity explanation must also reproduce the gravitational
potentials probed by CMB lensing, weak lensing, redshift-space distortions,
clusters, and galaxy dynamics.  In phenomenological language, the model must
provide the analogues of the effective Newton response and lensing response,
often denoted $\mu(a,k)$ and $\Sigma(a,k)$, together with any anisotropic stress,
sound speed, or extra-field propagation law.  These perturbative functions are
not determined by Eqs.~\eqref{eq:dark_energy_residual}--\eqref{eq:dark_pressure_residual};
they are additional model predictions.

Compactness adds a global consistency requirement.  Perturbations on the
unquotiented round $S^3$ use the discrete eigenvalues and kernels above, not a
continuous $\RR^3$ Fourier basis.  Nonlocal kernels, infrared prescriptions, and
vacuum choices in a modified-gravity model must therefore be defined on the
compact spectrum and must reduce continuously to the local flat-space
expressions as $R_0\to\infty$.

\section{Falsification criteria}

The compact reconstruction program fails, or becomes less useful than a simpler
parametrization, if any of the following occur:
\begin{enumerate}
\item Closed-FLRW distances are disfavored relative to flat or curved $\LCDM$
after equal-prior comparison.
\item The reconstructed $T^\D_{\mu\nu}$ violates conservation without an
explicit interaction term.
\item Perturbative data do not admit a stable pressureless clustering component
within the stated model class.
\item The inferred smooth component is inconsistent across CMB, BAO,
supernovae, lensing, and growth.
\item Compact $S^3$ covariance worsens posterior predictive tests.
\item Apparent dark-energy evolution disappears after curvature, neutrino mass,
calibration, and selection-function marginalization.
\item Visible-sector closure is inconsistent with BBN, CMB thermodynamics, or
neutrino priors.
\item A modified-gravity interpretation cannot supply an effective tensor and
perturbation sector matching the reconstructed residual, lensing, and growth
observables.
\item The local large-radius limit does not reduce continuously to the standard
FLRW expressions.
\end{enumerate}

\section{Conclusion}

The strict compact-geometric result is
\begin{equation}
\RR\times S^3
\Longrightarrow
\begin{gathered}
\text{finite closed geometry}\\
\text{and compact spectral covariance}.
\end{gathered}
\end{equation}
The strict GR residual result is
\begin{equation}
G_{\mu\nu}-\frac{8\pi G}{c^4}T^{\rm vis}_{\mu\nu}
\Longrightarrow
\begin{gathered}
T^\D_{\mu\nu}\ \text{as a covariantly}\\
\text{conserved residual}.
\end{gathered}
\end{equation}
The split into clustering matter and smooth acceleration is not a consequence
of topology.  It is obtained, within GR, only by combining homogeneous
distances with perturbative gravitational observables.  Compact topology is
therefore kept as exact geometry, while the dark sector is reconstructed as an
invariant residual and tested against public likelihoods.  The strongest
model-independent closure is the residual tensor itself; any further closure is
a dynamical model choice constrained by perturbative observables.

\begin{acknowledgments}
No external funding is declared.  The construction uses public cosmological
literature and standard differential geometry, general relativity, and
observational-cosmology methods.
\end{acknowledgments}

\section*{Data and code availability}
No proprietary data are used.  The paper defines an analytic and statistical
reconstruction pipeline.  A numerical implementation should use public CMB,
BAO, supernova, lensing, growth, cluster, and galaxy-dynamics likelihoods or
data products with stated priors.

\section*{Author contributions}
Andrei Preece and Boris Batenin contributed to the conceptual formulation,
invariant reconstruction framework, proof structure, and manuscript preparation.

\appendix

\section{Pressure reconstruction}

From Eq.~\eqref{eq:etot},
\[
\epsilon_{\rm tot}=\crit\left(E^2-\Omk a^{-2}\right).
\]
Therefore
\[
\frac{1}{\crit}\frac{\dd\epsilon_{\rm tot}}{\dd\ln a}
=\frac{\dd E^2}{\dd\ln a}+2\Omk a^{-2}.
\]
Using
\[
p_{\rm tot}=-\epsilon_{\rm tot}
-\frac13\frac{\dd\epsilon_{\rm tot}}{\dd\ln a}
\]
gives Eq.~\eqref{eq:ptot}.

\section{Closed distances and the local flat-metric limit}

The closed transverse comoving distance is
\[
D_M=R_0\sin(D_C/R_0).
\]
Since $\sin x=x-x^3/6+O(x^5)$, the large-radius limit $R_0\to\infty$ gives
$D_M\to D_C$ with corrections of order $(D_C/R_0)^2\sim|\Omega_K|$.  This is a
local geometric limit of the distance formula, not a topological identification
of compact $S^3$ with $\RR^3$.


\section{Flat-space limit of the invariant closure}
\label{app:flat_limit}

The large-radius limit used in this paper is local.  One sends
\[
R_0\to\infty,\qquad \Omega_K=-\frac{c^2}{H_0^2R_0^2}\to0,
\]
while holding fixed physical distances such as
$r=R_0\chi$, the line-of-sight comoving distance $D_C(z)$, and the continuum
wavenumber
\[
k=\frac{n+1}{R_0}.
\]
This is not a topological identification of $S^3$ with $\RR^3$; it is the
local geometric and spectral limit on scales much smaller than $R_0$.

The global geometric invariants behave as
\begin{align}
\Vol(S^3_{R_0})&=2\pi^2R_0^3\to\infty,
\nonumber\\
{}^{(3)}\mathcal R&=\frac{6}{R_0^2}\to0,
&
\diam(S^3_{R_0})&=\pi R_0\to\infty .
\end{align}
For a ball of fixed physical radius $r=R_0\chi$,
\[
V(\chi)=\frac{4\pi}{3}r^3+O(r^5/R_0^2),\qquad
A(\chi)=4\pi r^2+O(r^4/R_0^2).
\]
The causal-accessibility fraction in Eq.~\eqref{eq:access} has a subtle but
important limit.  For finite propagation time $T$,
\[
\Gamma(T)=\frac{\theta-\frac12\sin2\theta}{\pi}
=\frac{2}{3\pi}\left(\frac{cT}{R_0}\right)^3
+O(R_0^{-5})\to0,
\]
because the total spatial volume diverges.  The local accessible volume itself,
however, tends to the Euclidean value $4\pi(cT)^3/3$.  Similarly, the antipodal
condition in Eq.~\eqref{eq:conformal_access} is not automatically satisfied in
the flat limit; rather, the antipode recedes to infinite distance and the
condition ceases to be a finite-distance observable.

The distance map reduces continuously to the flat-FLRW one:
\begin{align}
D_M(z)&=R_0\sin\!\left(\frac{D_C(z)}{R_0}\right) \nonumber\\
&=D_C(z)\left[
1-\frac16\left(\frac{D_C(z)}{R_0}\right)^2
+O\!\left(\frac{D_C^4}{R_0^4}\right)\right],
\label{eq:flat_distance_expansion_full}
\end{align}
so that
\[
D_M\to D_C,\qquad
D_A\to\frac{D_C}{1+z},\qquad
D_L\to(1+z)D_C ,
\]
with relative corrections
\[
\left(\frac{D_C}{R_0}\right)^2
=(-\Omega_K)\left(\frac{H_0D_C}{c}\right)^2 .
\]
The BAO combinations therefore become
\[
D_V\to\left[zD_C^2D_H\right]^{1/3},\qquad
F_{\rm AP}\to\frac{D_CH}{c},
\]
while Etherington reciprocity is unchanged.

The background reconstruction reduces to the standard flat residual:
\begin{align}
\frac{\epsilon_{\rm tot}}{\crit}
&=E^2-\Omega_K a^{-2}\to E^2,\\
\frac{p_{\rm tot}}{\crit}
&=-\frac13\left[
\frac{\dd E^2}{\dd\ln a}+3E^2-\Omega_Ka^{-2}\right]
\nonumber\\
&\quad\to
-\frac13\left[
\frac{\dd E^2}{\dd\ln a}+3E^2\right],
\end{align}
and therefore
\[
\frac{\epsilon_{\D}}{\crit}
\to
E^2-\Ob a^{-3}-\Og a^{-4}-\Onu(a),
\]
with the corresponding flat expression for $p_{\D}$ obtained by subtracting
the visible pressure.  The acceleration scalar, the cold-clustering split, the
smooth residual, and the filters of Sec.~\ref{sec:complete_invariant_closure} reduce
with the same substitution $\Omega_K\to0$.

The scalar, vector, and tensor spectra also have a continuous local limit.  With
$k=(n+1)/R_0$,
\begin{align}
\frac{n(n+2)}{R_0^2}
&=k^2-\frac{1}{R_0^2}\to k^2,\\
\frac{n(n+2)-1}{R_0^2}
&=k^2-\frac{2}{R_0^2}\to k^2,\\
\frac{n(n+2)-2}{R_0^2}
&=k^2-\frac{3}{R_0^2}\to k^2 .
\end{align}
For scalar modes the compact replacement rule becomes
\begin{align}
&\frac{1}{\Vol(S^3_{R_0})}
\sum_{n=0}^\infty (n+1)^2
F\!\left(\frac{n+1}{R_0}\right)
\nonumber\\
&\quad\longrightarrow
\frac{1}{2\pi^2}\int_0^\infty k^2F(k)\,\dd k
=
\int\frac{\dd^3k}{(2\pi)^3}F(k).
\label{eq:flat_replacement_limit}
\end{align}
with the vector and tensor sectors giving the corresponding continuum
polarization sums.  Equivalently, the compact spectral covariance tends to the
flat Fourier covariance after the sum over hyperspherical shells is converted
to the integral over $k$; the familiar angular kernel $\sin(kr)/(kr)$ appears
inside the limiting integral,
\[
C(r)\to
\frac{1}{2\pi^2}\int_0^\infty
k^2 P(k)\frac{\sin(kr)}{kr}\,\dd k .
\]
Thus Eq.~\eqref{eq:closed_cmb_sum} becomes the usual flat line-of-sight
expression, up to the same normalization convention for transfer functions and
primordial spectra used in the closed calculation.

Finally, the invariant closure itself is continuous.  The order-by-order tensor
identity, the complete target $\mathcal C_{\rm inv}$, the posterior
factorization, and the dynamical filters do not acquire singular terms as
$\Omega_K\to0$.  Compact $S^3$ therefore supplies a finite closed infrared
bookkeeping system whose local observables and residual identities reduce
smoothly to the standard spatially flat FLRW framework.

\section{Requirements for a physical explanation}

A proposed physical origin for the residual is not part of this formalism until
it supplies:
\begin{enumerate}
\item an action, modified field equation, or stress tensor;
\item conservation laws or an explicit interaction term;
\item background functions $\epsilon(a)$ and $p(a)$;
\item perturbation equations, sound speed, and anisotropic stress;
\item lensing and growth predictions;
\item consistency with BBN, CMB, BAO, supernovae, clusters, and galaxies;
\item a local flat-metric and closed-geometry comparison;
\item public likelihood results against $\Lambda$CDM, curved $\Lambda$CDM,
$w$CDM, modified gravity, and nonparametric dark-energy reconstructions.
\end{enumerate}

\sloppy
\begin{thebibliography}{99}

\bibitem{Wald1984}
R. M. Wald, \textit{General Relativity} (University of Chicago Press, Chicago, 1984).

\bibitem{HawkingEllis1973}
S. W. Hawking and G. F. R. Ellis, \textit{The Large Scale Structure of Space-Time}
(Cambridge University Press, Cambridge, 1973).

\bibitem{doCarmo1992}
M. P. do Carmo, \textit{Riemannian Geometry} (Birkh\"auser, Boston, 1992).

\bibitem{Peebles1993}
P. J. E. Peebles, \textit{Principles of Physical Cosmology} (Princeton University Press,
Princeton, 1993).

\bibitem{DodelsonSchmidt2020}
S. Dodelson and F. Schmidt, \textit{Modern Cosmology}, 2nd ed. (Academic Press, London, 2020).

\bibitem{Mukhanov2005}
V. Mukhanov, \textit{Physical Foundations of Cosmology} (Cambridge University Press,
Cambridge, 2005).

\bibitem{Weinberg2008}
S. Weinberg, \textit{Cosmology} (Oxford University Press, Oxford, 2008).

\bibitem{RubinOrdonez1984}
M. A. Rubin and C. R. Ordonez, Symmetric tensor eigenspectrum of the Laplacian on
$n$-spheres, J. Math. Phys. \textbf{25}, 2888 (1984).


\bibitem{Higuchi1987}
A. Higuchi, Symmetric tensor spherical harmonics on the $N$-sphere and their application
to the de Sitter group SO($N,1$), J. Math. Phys. \textbf{28}, 1553 (1987);
Erratum J. Math. Phys. \textbf{43}, 6385 (2002).

\bibitem{AbbottSchaefer1986}
L. F. Abbott and R. K. Schaefer, A general, gauge-invariant analysis of the cosmic
microwave anisotropy, Astrophys. J. \textbf{308}, 546 (1986).

\bibitem{Jantzen1978}
R. T. Jantzen, Tensor harmonics on the 3-sphere, J. Math. Phys. \textbf{19}, 1163 (1978).

\bibitem{LindblomTaylorZhang2017}
L. Lindblom, N. W. Taylor, and F. Zhang, Scalar, vector and tensor harmonics on the
three-sphere, Gen. Relativ. Gravit. \textbf{49}, 139 (2017), arXiv:1709.08020.

\bibitem{Efstathiou2003}
G. Efstathiou, Is the low cosmic microwave background quadrupole a signature of spatial
curvature?, Mon. Not. R. Astron. Soc. \textbf{343}, L95 (2003), arXiv:astro-ph/0303127.

\bibitem{Planck2018VI}
N. Aghanim \textit{et al.} (Planck Collaboration), Planck 2018 results. VI. Cosmological
parameters, Astron. Astrophys. \textbf{641}, A6 (2020), arXiv:1807.06209.

\bibitem{Planck2018VII}
Y. Akrami \textit{et al.} (Planck Collaboration), Planck 2018 results. VII. Isotropy and
statistics of the CMB, Astron. Astrophys. \textbf{641}, A7 (2020), arXiv:1906.02552.

\bibitem{DESIDR2Lya2025}
M. Abdul-Karim \textit{et al.} (DESI Collaboration), DESI DR2 results. I. Baryon acoustic
oscillations from the Lyman alpha forest, Phys. Rev. D \textbf{112}, 083514 (2025),
arXiv:2503.14739.

\bibitem{DESIDR2BAO2025}
M. Abdul-Karim \textit{et al.} (DESI Collaboration), DESI DR2 results. II. Measurements of
baryon acoustic oscillations and cosmological constraints, Phys. Rev. D \textbf{112},
083515 (2025), arXiv:2503.14738.

\bibitem{DESIDR2DE2025}
K. Lodha \textit{et al.} (DESI Collaboration), Extended dark energy analysis using DESI
DR2 BAO measurements, Phys. Rev. D \textbf{112}, 083511 (2025), arXiv:2503.14743.

\bibitem{PantheonPlus2022}
D. Brout \textit{et al.}, The Pantheon+ analysis: Cosmological constraints,
Astrophys. J. \textbf{938}, 110 (2022), arXiv:2202.04077.

\bibitem{PantheonPlusData2022}
D. Scolnic \textit{et al.}, The Pantheon+ analysis: The full data set and light-curve
release, Astrophys. J. \textbf{938}, 113 (2022), arXiv:2112.03863.

\bibitem{Moresco2022}
M. Moresco \textit{et al.}, Unveiling the Universe with emerging cosmological probes,
Living Rev. Relativ. \textbf{25}, 6 (2022), arXiv:2201.07241.

\bibitem{ACTDR62024}
M. S. Madhavacheril \textit{et al.} (ACT Collaboration), The Atacama Cosmology Telescope:
DR6 gravitational lensing map and cosmological parameters, Astrophys. J. \textbf{962},
113 (2024), arXiv:2304.05203.

\bibitem{DESY32022}
T. M. C. Abbott \textit{et al.} (DES Collaboration), Dark Energy Survey Year 3 results:
Cosmological constraints from galaxy clustering and weak lensing, Phys. Rev. D
\textbf{105}, 023520 (2022), arXiv:2105.13549.

\bibitem{BulletCluster2006}
D. Clowe, M. Brada\v{c}, A. H. Gonzalez, M. Markevitch, S. W. Randall, C. Jones, and
D. Zaritsky, A direct empirical proof of the existence of dark matter, Astrophys. J. Lett.
\textbf{648}, L109 (2006), arXiv:astro-ph/0608407.

\bibitem{SPARC2016}
F. Lelli, S. S. McGaugh, and J. M. Schombert, SPARC: Mass models for 175 disk galaxies
with Spitzer photometry and accurate rotation curves, Astron. J. \textbf{152}, 157
(2016), arXiv:1606.09251.

\bibitem{Cooke2018}
R. J. Cooke, M. Pettini, and C. C. Steidel, One percent determination of the primordial
deuterium abundance, Astrophys. J. \textbf{855}, 102 (2018), arXiv:1710.11129.

\bibitem{SchrammTurner1998}
D. N. Schramm and M. S. Turner, Big-bang nucleosynthesis enters the precision era,
Rev. Mod. Phys. \textbf{70}, 303 (1998), arXiv:astro-ph/9706069.

\bibitem{Etherington1933}
I. M. H. Etherington, On the definition of distance in general relativity,
Philos. Mag. \textbf{15}, 761 (1933).

\bibitem{Fixsen2009}
D. J. Fixsen, The temperature of the cosmic microwave background, Astrophys. J.
\textbf{707}, 916 (2009), arXiv:0911.1955.

\bibitem{Mangano2005}
G. Mangano, G. Miele, S. Pastor, T. Pinto, O. Pisanti, and P. D. Serpico,
Relic neutrino decoupling including flavor oscillations, Nucl. Phys. B \textbf{729},
221 (2005), arXiv:hep-ph/0506164.

\bibitem{deSalasPastor2016}
P. F. de Salas and S. Pastor, Relic neutrino decoupling with flavor oscillations
revisited, JCAP \textbf{07}, 051 (2016), arXiv:1606.06986.

\bibitem{Froustey2020}
J. Froustey, C. Pitrou, and M. C. Volpe, Neutrino decoupling including flavor
oscillations and primordial nucleosynthesis, JCAP \textbf{12}, 015 (2020),
arXiv:2008.01074.

\bibitem{LesgourguesPastor2006}
J. Lesgourgues and S. Pastor, Massive neutrinos and cosmology, Phys. Rep. \textbf{429},
307 (2006), arXiv:astro-ph/0603494.

\bibitem{LesgourguesPastor2012}
J. Lesgourgues and S. Pastor, Neutrino mass from cosmology, Adv. High Energy Phys.
\textbf{2012}, 608515 (2012), arXiv:1212.6154.

\bibitem{HuSawicki2007}
W. Hu and I. Sawicki, Models of $f(R)$ cosmic acceleration that evade solar-system tests,
Phys. Rev. D \textbf{76}, 104043 (2007), arXiv:0708.1190.

\bibitem{Linder2003}
E. V. Linder, Exploring the expansion history of the Universe, Phys. Rev. Lett.
\textbf{90}, 091301 (2003), arXiv:astro-ph/0208512.

\bibitem{Lewis2011}
A. Lewis, A. Challinor, and A. Lasenby, Efficient computation of CMB anisotropies in
closed FRW models, Astrophys. J. \textbf{538}, 473 (2000),
arXiv:astro-ph/9911177.

\bibitem{Lesgourgues2011}
J. Lesgourgues, The Cosmic Linear Anisotropy Solving System (CLASS) I:
Overview, arXiv:1104.2932 (2011); D. Blas, J. Lesgourgues, and T. Tram,
The Cosmic Linear Anisotropy Solving System (CLASS) II: Approximation schemes,
JCAP \textbf{07}, 034 (2011), arXiv:1104.2933.

\end{thebibliography}

\end{document}
